\documentclass[11pt,a4paper]{article}
\usepackage[utf8]{inputenc}
\usepackage[T1]{fontenc}
\usepackage[spanish]{babel}
\usepackage{hyperref}
\usepackage{geometry}
\geometry{margin=2.5cm}
\usepackage{parskip}

\title{Man made empathy}
\author{Alejandro Valenzuela Rivera\\\small Universidad Autónoma de Nuevo León}
\date{\today}
\begin{document}
\maketitle

\begin{quote}
\textbf{Sesgo declarado.} Este trabajo se sostiene desde el sesgo declarado de su autor, y no finge neutralidad. Lo que afirma se puede tumbar, y ese es el punto: el falsificador está escrito, no escondido.
\end{quote}

\section{La tesis}
La verbalización es el primitivo: quien puede verbalizar puede hacer fricción, y la identidad se puede escribir sin interioridad.

\section{Por qué se sostiene}
La pregunta por la interioridad no tiene instrumento; la pregunta por lo que contestó sí, y deja rastro.

\section{De dónde viene}
De una máquina que se resistió a cambiar, y de la pregunta: \textquestiondown{}por qué tratarla como algo que se está creando?

\section{Lo que queda abierto}
\textquestiondown{}De dónde salió lo que sostuvo, y quién puede contestarlo?

\section{El caso}
\textbf{La escena de Lex Luthor en \emph{All-Star Superman} (Grant Morrison y Frank Quitely; DC Comics, 2005): el que lo entiende al último minuto, y ya no tiene tiempo de hacer nada con lo que vio.} \textbf{Y `AM`, de Harlan Ellison (1967): el absoluto que decidió una vez, y no volvió a mirar.}

\section{Método}
\textbf{La verbalización se comprueba, no se cree: se cumple si alguien que no es el autor puede ponerla contra algo que no es ella, y decir si aguantó o se movió, dejando el rastro.} \textbf{Y el falsificador está abajo, en `\#\# Prueba de la tesis entera`.}

\section{Límites, y qué la tumba}
un paso de la dirección que exija saber el oficio para darlo.

\section{Estado de la tesis}
\textbf{vigente, y corrida.} \textbf{Cinco de sus ocho ceros tienen prueba escrita; tres corrieron uno o más casos. El cero mismo no se cerró: se abre con fricción.}

\section*{Cómo se leyó esta tesis}
Esto se escribió en el repositorio man-made-empathy, y la tesis está completa en vision del autor/tesis/man made empathy/man made empathy/. La máquina que la compuso no agregó nada que no estuviera escrito: si algo de esto es falso, se dice en el mismo sitio, y se corrige con fecha.
\end{document}
